\paragraph{The analytic enclosures and their actual matrix maps.}
For positive weights $w_i$ let $W=\operatorname{diag}(w_i)$. The
specified norm is the maximum absolute row-sum norm of the actual
conjugate $W^{-1}MW$:
\[
 \|M\|_w=\max_i\sum_j|M_{ij}|\frac{w_j}{w_i}.
\]
The equality $W^{-1}MNW=(W^{-1}MW)(W^{-1}NW)$ and the rowwise
triangle inequality prove
$\|MN\|_w\le\|M\|_w\|N\|_w$ and
$|M_{ij}|\le\|M\|_w w_i/w_j$. For matrices with centres $C,D$
and entrywise error bounds $E,F$, multiplication has entrywise
error at most
\[
 |C|F+E|D|+EF,
\]
with the computed rounding error also retained. In particular, a
left factor $H_c$ acts on an entrywise error matrix by the actual
sum $|H_c|_+E$, where $|H_c|_+$ bounds its entries including their
ball radii.

\paragraph{Root series.}
The listed regular-root connection matrices already include the
specified diagonal gauges. Their frame map is
$E_{\rm root}=E_{\rm original}D_{\rm root}$, and the numerical
coordinate change is exactly
$G_z(z)=\alpha G_{\rm root}(\alpha z)$, where
$\alpha^4=256/3125$ at zero and $\alpha^5=3125/256$ at infinity.
Multiplying $G_{\rm root}(t)$ by $1+t^m/\alpha^m$ gives the exact
rational polynomial $\sum_{j=0}^{m-1}C_jt^j$, so
$B_j=\alpha^{j+1}C_j$. Thus
\[
 U'=U\frac{\sum_{j=0}^{m-1}B_jz^j}{1+z^m},\qquad U(0)=I.
\]
Coefficient comparison in $(1+z^m)U'=UB$ gives
\[
 (n+1)U_{n+1}
 =\sum_{j=0}^{\min(m-1,n)}U_{n-j}B_j
   -(n-m+1)U_{n-m+1},
\]
where the last term is absent for $n<m$. The order of every
matrix product is as displayed. Picard iteration on any smaller
disc converges by the bound $(Lr)^k/k!$ on its $k$-fold integral,
so these are the coefficients of the actual normalized solution.

Use $w=(1,1/128,1/128,1/32,1/64,1/32)$ and
\[
 \mathfrak b=\max_i\sum_{j,k}|(B_k)_{ij}|\frac{w_j}{w_i}.
\]
On a ray of radius $t<1$ the connection norm is at most
$\mathfrak b/(1-t^m)$. The scalar Gronwall inequality
\cite[Lemma~2.7, pp.~42--43]{TeschlODE2012} applies to the actual
matrix integral equation. The elementary estimate
\[
 \frac1{1-t^m}\le\frac1{m(1-t)}+\frac{m-1}{m}
\]
holds because its numerator difference is
$\bigl(\sum_{j=1}^{m-1}t^j-(m-1)t^m\bigr)/m\ge0$.
Consequently
\[
 \|U(z)\|_w\le K(R):=
 \exp\!\left\{\mathfrak b\left[-\frac{\log(1-R)}m
                 +\frac{m-1}{m}R\right]\right\}
 \quad (|z|\le R<1).
\]
Cauchy's coefficient integral, with the positively oriented circle
of radius $R$ \cite[\S1.9(iii), (1.9.30)--(1.9.31)]{DLMFCauchy2026},
gives $\|U_n\|_w\le K(R)R^{-n}$. Hence the error after degree $N$
on the whole disc $|z|\le r<R$ is
\[
 \tau_N(r;R)=K(R)\frac{(r/R)^{N+1}}{1-r/R}.
\]
Its $ij$ entry is bounded by $\tau_Nw_i/w_j$; after left
multiplication by $H_c$ the bound is
$\tau_N\sum_l|(H_c)_{il}|_+w_l/w_j$.
The current root computations use $N=3000$, $R=0.999$ and each
rectangle's actual maximum radius. Every coefficient and point
rounding radius remains in the corresponding ball calculation.

\paragraph{Enclosing an entire rectangle.}
Let $z_c$ be its centre, $\rho$ its Euclidean half-diagonal, and
$\mathcal A$ a nonnegative entrywise bound for $|G(z)|$ on the
whole rectangle. For any point $z$ the segment
$z_c+t(z-z_c)$ remains in it. Iteration of the integral equation
on that segment bounds the $k$-th ordered integral by
$|Y(z_c)|_+|z-z_c|^k\mathcal A^k/k!$. Subtracting its zeroth term
and summing proves the entrywise simultaneous enclosure
\[
 |Y(z)-Y(z_c)|
 \le |Y(z_c)|_+\bigl(\exp(\rho\mathcal A)-I\bigr).
\]
The centre enclosure and matrix-exponential rounding radii are
retained, so this bounds every point of the rectangle.

The midpoint-polynomial bound has the same explicit domain
control. For $p(z)=\sum_{j=0}^Nc_jz^j$, coefficient discs with
midpoints $\widehat c_j$ and radii $\epsilon_j$, and midpoint
evaluation rounding radius $\eta$, the centre error is
$\eta+\sum_j\epsilon_j|z_c|^j$. Integrating $p'$ on each segment
bounds the remaining variation by
\[
 \rho\sum_{j=1}^N j\bigl(|\widehat c_j|+\epsilon_j\bigr)r^{j-1},
\]
where $r$ is the rectangle's maximum modulus. Add the analytic
tail at that same $r$. This formula includes coefficient
uncertainty in the derivative bound as well as at the centre.

\paragraph{The cusp holomorphic factor.}
The exact connection here is
$R/q+A_0/(1+q)$, with $R^2=0$. In
$Y=H(I+R\log q)P(q)$, $P(0)=I$, the inverse of the logarithmic
factor is $I-R\log q$. Multiplication by that inverse gives
\[
 P'=\frac{\operatorname{ad}(P)}q+\frac{PA_0}{1+q},
 \qquad\operatorname{ad}(Q)=QR-RQ.
\]
Thus
\[
 (nI-\operatorname{ad})P_n
  =\sum_{j=0}^{n-1}(-1)^{n-1-j}P_jA_0.
\]
Equivalently its right side is
$\operatorname{ad}(P_{n-1})+P_{n-1}A_0-(n-1)P_{n-1}$.
Direct expansion gives
$\operatorname{ad}^2(Q)=-2RQR$ and
$\operatorname{ad}^3(Q)=0$, so the actual inverse is
\[
 (nI-\operatorname{ad})^{-1}
  =\frac1n\left(I+\frac{\operatorname{ad}}n
                   +\frac{\operatorname{ad}^2}{n^2}\right).
\]
This fixes both the commutator sign and the side of the $A_0$
multiplication.

Use the actual cusp weights $(100,1,1,1,1,1)$ and put
$r_0=\|R\|_w$, $a_0=\|A_0\|_w$,
$C_0=1+2r_0+4r_0^2$, $\gamma=C_0a_0$. For the supplied exact
matrices,
\[
 r_0=\frac{5103}{8000},\quad a_0=\frac{367}{512},\quad
 C_0=\frac{62452609}{16000000},\quad
 \gamma=\frac{22920107503}{8192000000}.
\]
The norm of the inverse operator is at most $C_0/n$. Therefore
$\|P_n\|_w\le(\gamma/n)\sum_{j<n}\|P_j\|_w$.
The nonnegative scalar majorant $c_0=1$,
$c_n=(\gamma/n)\sum_{j<n}c_j$ satisfies
$c_n=(\gamma+n-1)c_{n-1}/n$ and has generating function
$(1-s)^{-\gamma}$. Its radius is at least one by the coefficient
ratio. Thus the actual matrix series converges normally on
$|q|<1$, solves the equation term by term, and obeys
\[
 \|P(q)\|_w\le K_0(s):=(1-s)^{-\gamma}\quad(|q|\le s<1).
\]
The coefficient integral and geometric sum give its degree-$N$
tail on $|q|\le r<s$ as
\[
 K_0(s)\frac{(r/s)^{N+1}}{1-r/s},
\]
with the exact entry factor $w_i/w_j$. The current cusp series
uses $N=600$ and $s=0.4$ for $r\le0.25$. On the inner disc
$r=0.01$ the enclosure retains all actual coefficient absolute
sums through degree600 and this tail. It bounds $P$ itself;
the subsequent logarithmic factor retains its specified branch
and its own actual matrix entries.

\paragraph{Collar continuation.}
At a centre $z_c$ put $t=z-z_c$ and expand
\[
 1+(z_c+t)^5=\sum_{j=0}^5q_jt^j,
 \quad q_j=\binom5jz_c^{5-j}+\delta_{j0},
 \qquad
 \widetilde B_j=\sum_{k=j}^4\binom kjB_kz_c^{k-j}.
\]
For $T(0)=I$ the exact transition coefficients satisfy
\[
 T_{n+1}=
 \frac{\displaystyle
   \sum_{j=0}^{\min(4,n)}T_{n-j}\widetilde B_j
  -\sum_{j=1}^{\min(5,n)}q_j(n-j+1)T_{n-j+1}}
 {(n+1)q_0}.
\]
Let $\zeta_j=e^{(2j+1)\pi i/5}$ and
$R_{\rm loc}=\tfrac12\min_j|z_c-\zeta_j|$. On the closed disc
of this radius,
\[
 |1+(z_c+t)^5|\ge D_{\rm lower}
   :=\prod_j\bigl(|z_c-\zeta_j|-R_{\rm loc}\bigr)>0.
\]
The numerator norm is at most
\[
 B_{\max}=\max_i\sum_{j,k}|(\widetilde B_k)_{ij}|
                  \frac{w_j}{w_i}R_{\rm loc}^k.
\]
For $M=B_{\max}/D_{\rm lower}$, radial Gronwall and Cauchy
give the exact degree-$N$ transition tail
\[
 e^{MR_{\rm loc}}
 \frac{(|\delta|/R_{\rm loc})^{N+1}}
      {1-|\delta|/R_{\rm loc}}.
\]
The current steps use $N=160$ and
$|\delta|/R_{\rm loc}<0.3$. The actual continuation order is
$Y(z_c+\delta)=Y(z_c)T_{z_c}(\delta)$, as follows by checking
the right ODE and its initial value. Every transition tail is
added entrywise before multiplying on the right, and the
entrywise product rule propagates all earlier errors.

\paragraph{The real implicit system and Hessian.}
Use the specified fixed four rows on the central square and
write the actual system as $Mv=b$, with
\[
 M_{i,*}=(\operatorname{Im}Y_{i1},\operatorname{Re}Y_{i1},
          \operatorname{Im}Y_{i2},\operatorname{Re}Y_{i2}),
 \qquad b_i=\operatorname{Re}Y_{i3}.
\]
For the top row define $c$ by the same four coordinate projections,
$d=\operatorname{Re}Y_{03}$ and $h=k(cv-d)$,
$k=96/(625\pi)$. The certified determinant enclosure excludes
zero on the entire central square and on its sixteen subsquares;
the matrix inverse is therefore the actual one there.
For $a,b\in\{x,y\}$, direct differentiation gives
\[
 \begin{split}
 v_a&=M^{-1}(b_a-M_av),\\
 v_{ab}&=M^{-1}(b_{ab}-M_{ab}v-M_av_b-M_bv_a),\\
 h_a&=k(c_av+cv_a-d_a),\\
 h_{ab}&=k(c_{ab}v+c_av_b+c_bv_a+cv_{ab}-d_{ab}).
 \end{split}
\]
The repeated-index cases retain both cross terms. The actual
complex derivatives supplying these real projections are
\[
 Y_x=YG,\quad Y_y=iYG,\quad
 Y_{xx}=Y(G^2+G'),\quad Y_{xy}=iY(G^2+G'),\quad
 Y_{yy}=-Y(G^2+G').
\]
Taking the stated real and imaginary parts and carrying out these
products with rigorous full-square balls consequently encloses
all Hessian entries throughout the domain. Approximate row
selection merely chooses a minor; it supplies neither its inverse
nor its error bound.

The complete domain certificate uses110-digit ball arithmetic and
the parameters above. Its interval values, subdivision counts and
strict comparison threshold are those of the separately recorded
current domain run. The independent analytic witness
\texttt{P1\_18\_analytic\_certificates\_fresh.py} uses256-bit balls
and713 checks to recompute both degree3000 root series, the
degree600 cusp series, representative degree160 collar
transitions and their actual compositions, the full commutator
operator and the real implicit derivative maps. Its representative
domains are recorded separately from the complete leaf partition.
